\section{Combinación de gradiente de política y Actor-Critic}
\subsection{Introducción a la arquitectura Actor-Critic}
La estructura Actor-Critic entrena en conjunto la política (actor) y la función de valor (critic). El critic proporciona una estimación de advantage y el actor ajusta la política según ese advantage, formando una estructura de dos torres:
\begin{enumerate}
    \item El actor produce la distribución de acciones $\pi_\theta(a\mid s)$.
    \item El critic estima $V_w(s)$, con lo que se calcula el advantage $\hat{A}_t = G_t - V_w(s_t)$.
    \item El actor se actualiza usando el advantage anterior, y el critic actualiza la red de valor con la pérdida TD.
\end{enumerate}
\begin{knowledgebox}{La sinergia del Actor-Critic}
El actor, en tanto política, aporta la decisión de alto nivel; el critic evalúa la recompensa estimada; ambos en conjunto mitigan la alta varianza del gradiente de política tradicional.
\end{knowledgebox}

La actualización con bootstrap del critic define el error TD con $\delta = r_t + \gamma V_w(s_{t+1}) - V_w(s_t)$, y minimizar $\delta^2$ permite que el critic ofrezca una línea base más precisa.

\subsection{Actor-Critic frente al gradiente de política original}
El gradiente de política original solo usa recompensa por Monte Carlo, por lo que tiene alta varianza y aprovecha mal los datos; el Actor-Critic introduce la estimación de valor como línea base, lo que reduce la varianza y permite el bootstrapping.
\begin{table}[H]
\centering
\small
\begin{tabular}{p{0.34\textwidth}p{0.34\textwidth}p{0.3\textwidth}}
\toprule
Método & Tipo de estimación & Aprovechamiento de muestras \\
\midrule
REINFORCE & Completamente Monte Carlo & Se usa cada trayectoria una vez \\
Actor-Critic & Valor con bootstrap & Los mismos datos se pueden actualizar repetidamente \\
PPO & Sustituto recortado + cabeza de valor & Aprovecha varias épocas \\
\bottomrule
\end{tabular}
\caption{Comparación de los métodos de la familia de gradiente de política}
\end{table}

En el entrenamiento real, la expresión del gradiente del actor sigue siendo $\mathbb{E}[\nabla_\theta \log \pi_\theta(a\mid s) \cdot \hat{A}]$, pero el advantage lo proporciona el critic, y puede usarse la forma GAE para mejorar la estabilidad. Al actualizar el critic, puede fusionarse con TD de varios pasos (como TD($\lambda$)) la señal de Monte Carlo con la de bootstrap, formando así sinergia dentro de la arquitectura de dos torres.

\subsection{Resumen de la sección}
El Actor-Critic combina el gradiente de política con la función de valor, reduce la varianza y mejora la eficiencia de muestras; PPO/TRPO agregan además una región de confianza sobre esta base, formando un ciclo de entrenamiento más estable.

